\section{Related Work}

\paragraph{Multi-agent LLM reasoning.}
A growing body of work improves LLM reasoning by letting several model instances critique one another.
Multi-agent debate lets agents exchange and revise answers over several rounds~\cite{du2024improving}; MAD assigns opposing roles to two debaters and adds a judge that can end the debate~\cite{liang2024mad}; ReConcile runs a round-table discussion among diverse LLMs and aggregates their answers by confidence-weighted voting~\cite{chen2024reconcile}.
Self-refinement and reflection methods apply the same critique--revise loop within a single model~\cite{madaan2023selfrefine,shinn2023reflexion}.
Role-based systems such as ChatDev, MetaGPT, and AutoGen organize agents into specialized roles that collaborate through structured conversations~\cite{qian2024chatdev,hong2024metagpt,wu2024autogen}.
Closest to our setting, MARS organizes agents as an author, independent reviewers, and a meta-reviewer, which removes reviewer-to-reviewer communication and thereby reduces the token cost of fully connected debate~\cite{wang2025mars}.
Our framework adopts this reviewer-style role structure and uses it as the message-based baseline.

\paragraph{Role specialization, workflow optimization, and graph-structured orchestration.}
Recent systems go beyond fixed role assignments.
AgentVerse dynamically recruits a group of expert agents for each task and alternates collaborative decision making with evaluation~\cite{chen2024agentverse}; DyLAN selects and re-weights agents during inference through an agent-importance score~\cite{liu2024dylan}.
Role-specialized pipelines have also been shown to be effective beyond question answering: Zhao et al.\ combine multiple LLM agents with small models for zero-shot knowledge-graph question generation, assigning distinct sub-tasks to distinct agents~\cite{zhao2024kgqg}.
A second line of work treats the communication structure itself as the object of optimization.
MacNet arranges agents on directed acyclic graphs and studies how collaboration scales with topology~\cite{qian2025scaling}; GPTSwarm represents an agent system as a computational graph whose node prompts and edge connectivity are optimized~\cite{zhuge2024gptswarm}; AgentPrune removes redundant messages from the spatial--temporal communication graph to reduce cost~\cite{zhang2025agentprune}.
These approaches decide \emph{who} talks to \emph{whom}.
Our work is complementary: we keep the topology fixed and ask \emph{what} the agents coordinate through---an explicit, validated task state rather than the raw message stream---and whether the degree of centralization of that state should depend on task uncertainty.
An uncertainty signal such as the CUI could in principle be used as an input to topology-selection methods.

\paragraph{Memory-based agents.}
Memory mechanisms also structure the context of LLM agents and are therefore closely related to our use of explicit state.
Generative Agents store a stream of observations and retrieve from it by recency, relevance, and importance~\cite{park2023generative}; MemGPT pages information between the context window and external storage in the manner of an operating system~\cite{packer2023memgpt}; MemoryBank maintains and updates a long-term store with a forgetting mechanism~\cite{zhong2024memorybank}; Reflexion and Voyager accumulate verbal lessons and executable skills across episodes~\cite{shinn2023reflexion,wang2024voyager}; and CoALA provides a unifying taxonomy of working, episodic, semantic, and procedural memory for language agents~\cite{sumers2024coala}.
In CoALA's terms, our state store is a structured \emph{working memory} for a single task, and we do not claim a new memory architecture.
The difference lies in what the structure is used for.
Memory systems answer the question ``which past information should be placed in the context?''; their central operations are write, retrieve, and summarize, and their goal is to preserve information across long horizons.
The world model in our framework answers the question ``may this proposed step be committed, and may the current trajectory be accepted?''; its central operations are a typed transition, a validator that can reject an action before other agents see it, and a risk function over the trajectory that gates acceptance independently of any single agent's verdict.
Our state is also not persistent: it is created per query and discarded afterwards, so none of the reported gains can come from cross-task memory.
The two ideas are compatible, and combining a persistent memory with state-level validation is a natural extension.

\paragraph{World models and structured reasoning.}
World models were introduced in reinforcement learning to represent environment states and transitions and to plan by imagined rollouts~\cite{ha2018recurrent,hafner2020dream}.
In the LLM literature, RAP uses the LLM itself as a world model that predicts the next reasoning state and plans with Monte-Carlo tree search~\cite{hao2023rap}, Tree-of-Thoughts evaluates and searches over partial solutions~\cite{yao2023tree}, and ReAct and Toolformer couple reasoning with environment or tool feedback~\cite{yao2023react,schick2023toolformer}.
These methods structure the reasoning of a single agent.
Our world model is far simpler than a learned dynamics model---it is a symbolic state store with rule-based validation and risk estimation---but it is shared across agents and used for coordination: it determines what reviewers see and when a candidate is accepted.

\paragraph{Positioning.}
Table~\ref{tab:related_work} summarizes the comparison.
To our knowledge, prior multi-agent reasoning frameworks do not combine an explicit per-task state object, validation of state transitions, a trajectory-level acceptance test, and a task-dependent choice of coordination mechanism.
The table also makes the limits of our framework explicit: it has neither persistent cross-task memory nor a learned communication topology.

\begin{table*}[!t]
\caption{Positioning Against Representative Approaches (\cmark: supported; \xmark: not supported; $\circ$: partially)}
\label{tab:related_work}
\centering
\footnotesize
\setlength{\tabcolsep}{2.5pt}
\renewcommand{\arraystretch}{1.15}
\begin{tabular}{lccccccc}
\toprule
 & \makecell{Debate / MAD\\\cite{du2024improving,liang2024mad}} & \makecell{ReConcile\\\cite{chen2024reconcile}} & \makecell{MARS\\\cite{wang2025mars}} & \makecell{AgentVerse\\\cite{chen2024agentverse}} & \makecell{GPTSwarm\\\cite{zhuge2024gptswarm}} & \makecell{Memory agents\\\cite{packer2023memgpt,zhong2024memorybank}} & Ours \\
\midrule
Multi-agent critique of candidate answers & \cmark & \cmark & \cmark & \cmark & \cmark & \xmark & \cmark \\
Explicit per-task state object & \xmark & \xmark & \xmark & \xmark & \xmark & $\circ$ & \cmark \\
Validation of state transitions & \xmark & \xmark & \xmark & \xmark & \xmark & \xmark & \cmark \\
Acceptance independent of a single agent's verdict & $\circ$ & \cmark & \xmark & $\circ$ & $\circ$ & \xmark & \cmark \\
Persistent cross-task memory & \xmark & \xmark & \xmark & \xmark & \xmark & \cmark & \xmark \\
Dynamic or optimized agent topology & \xmark & \xmark & \xmark & \cmark & \cmark & \xmark & \xmark \\
Uncertainty-dependent coordination mode & \xmark & \xmark & \xmark & \xmark & \xmark & \xmark & \cmark \\
\bottomrule
\end{tabular}
\end{table*}
